\chapter{神经元打乒乓}
% 完整版：chapters/22_dishbrain.tex

芯片上活神经元最著名的实验是 DishBrain，即“培养皿里的大脑”。大约一百万个神经元长在一块有两万六千个电极的芯片上，玩最简单的电脑乒乓游戏：球在屏幕上飞来飞去，球拍上下移动。

\section*{游戏是怎样设置的}

球通过八个电极进入网络。哪个电极在咔嗒——说明球在什么高度。咔嗒得有多频繁——说明球有多远：球在远端墙边时每秒四次，飞向球拍时越来越频繁，最多到四十次。输出是两片电极区域：一片里的咔嗒让球拍向上移动，另一片里的让它向下。计算机每十毫秒数一次咔嗒并移动球拍。这十毫秒是试验台计算机的时钟节拍，而不是培养物的：网络本身和所有神经网络一样，是异步工作的。

\pencilfig{22}{\pencilart{22}}{球、球拍，还有推动球拍的一百万个神经元。}

\section*{老师}

没有任何多巴胺。成功击中之后，网络收到一个短暂而完全可预测的信号：八个电极同时刺激，持续十分之一秒。未击中之后——是在随机位置、随机时刻进行的四秒杂乱刺激，然后停顿，再重新发球。

\section*{结果如何}

在二十分钟的一轮实验中，活培养物的回合变长了，漏接的球变少了。有完整反馈时效果最好。如果在击中和未击中之后不给信号、只是一片安静，也有改善，但弱一些；没有反馈时则没有改善。人类神经元比小鼠神经元玩得稍好一点。不过，延续好几天的学习被证明并不可靠。

\section*{网络为什么会学}

作者们这样解释结果：网络力求让自己的感觉变得可预测。未击中带来混乱，击中带来秩序，网络为了让秩序更多而改变自己。但还有一个更朴素的解释，也与数据相符：“没击中就摇一摇，击中了就别动”。每次摇晃都会把连接随机打乱一点。让网络频频失手的设置不断被打乱，而成功的设置则原封不动。网络自然而然地停留在失手少的地方——没有谁教它，只是不再有人摇它了。在我们的简单模型里，这样一条规则确实提高了击中率。

\pencilfig{130}{%
% scrambler: miss probability p over one setting of the network (valley in the middle); a miss kicks the setting
% at random, a hit leaves it alone, so the time spent at a setting is proportional to 1/p (6x more in the valley here)
\draw[pen] (0,1.2) -- (8.2,1.2);
\draw[penlite=pcRed] plot[smooth] coordinates {(0.0,2.46) (0.8,2.46) (1.6,2.46) (2.0,2.44) (2.4,2.41) (2.8,2.32) (3.2,2.15) (3.6,1.90) (4.0,1.62) (4.4,1.44) (4.8,1.44) (5.2,1.62) (5.6,1.90) (6.0,2.15) (6.4,2.32) (6.8,2.41) (7.2,2.44) (7.6,2.46) (8.0,2.46)};
\node[lbl, text=pcRed, anchor=south] at (7.55,2.55) {失手多};
\node[lbl, text=pcRed, anchor=west] at (5.3,1.45) {失手少};
% kicked balls on the slopes
\draw[dotpen=pcBlue, wash=pcBlue] (1.3,2.68) circle (0.12);
\draw[arr=pcGray, densely dashed] (1.35,2.82) to[bend left=50] (2.35,2.72);
\draw[arr=pcGray, densely dashed] (1.22,2.82) to[bend right=50] (0.4,2.72);
\draw[dotpen=pcBlue, wash=pcBlue] (6.3,2.45) circle (0.12);
\draw[arr=pcGray, densely dashed] (6.25,2.59) to[bend right=50] (5.45,2.25);
\node[lbl, anchor=south] at (1.3,3.05) {失手后摇一摇};
% the ball left alone in the valley
\draw[dotpen=pcBlue, wash=pcBlue] (4.6,1.66) circle (0.12);
\node[lbl, anchor=south] at (4.6,1.95) {不摇};
% where the network spends its time (proportional to 1/p)
\fill[pcGreen, opacity=0.22] (0,0) -- plot[smooth] coordinates {(0.0,0.18) (0.8,0.18) (1.6,0.18) (2.0,0.19) (2.4,0.19) (2.8,0.21) (3.2,0.24) (3.6,0.33) (4.0,0.55) (4.4,0.98) (4.8,0.98) (5.2,0.55) (5.6,0.33) (6.0,0.24) (6.4,0.21) (6.8,0.19) (7.2,0.19) (7.6,0.18) (8.0,0.18)} -- (8.0,0) -- cycle;
\draw[penlite=pcGreen] plot[smooth] coordinates {(0.0,0.18) (0.8,0.18) (1.6,0.18) (2.0,0.19) (2.4,0.19) (2.8,0.21) (3.2,0.24) (3.6,0.33) (4.0,0.55) (4.4,0.98) (4.8,0.98) (5.2,0.55) (5.6,0.33) (6.0,0.24) (6.4,0.21) (6.8,0.19) (7.2,0.19) (7.6,0.18) (8.0,0.18)};
\draw[pen] (0,0) -- (8.2,0);
\node[lbl, text=pcGreen!60!black, anchor=west] at (5.3,0.6) {网络停留的地方};
\node[lbl, anchor=north] at (4.1,-0.03) {网络的设置};
}{未击中会随机打乱网络的设置，击中则让它们保持原样。失手多的地方，网络待不住。失手少的地方，不再有人摇它，它便留在那里——尽管没有谁把它引到那里。}

怎样区分这两种解释？需要一个还没人做过的实验：未击中之后，给的不是杂乱的信号，而是时长和强度都相同的\emph{可预测}信号。如果网络照样能学会，那么关键就不在可预测性，而只在于“击中之后”和“未击中之后”的差别。这个实验还在等待有人去做。

我们的旅程到此结束。这台二十瓦的机器是用简单的零件搭成的，但它究竟怎样把它们组合成思想，我们目前只懂得一部分。也许，下一章会由你来写。

\fullversion{22}
